% sec_method.tex — Vitrine v5: Method
% Capture-adaptive video→3D→UE 5.8 pipeline.
% Companion to sec_results.tex. Included by the v5 master document.
%
% Cite keys: from report/references.bib + v5 additions noted in-text.
% Figure labels: \ref{fig:router}, \ref{fig:nanogs_room}, \ref{fig:scene_integrated},
%                \ref{fig:mesh_comparison}, \ref{fig:frame_gate}.
% British English throughout.

\section{Overview: the Capture-Adaptive Thesis}
\label{sec:method:overview}

The fundamental claim of Vitrine is that \emph{there is no single best reconstruction
pipeline for real-world video}.  Real captures fail in qualitatively different ways ---
motion blur, insufficient coverage, missing depth sensing, hardware constraints --- and
any single fixed recipe will be sub-optimal for most of them.  Vitrine therefore does not
select one pipeline at design time.  Instead, every incoming capture is \emph{diagnosed}
first, and a \emph{router} selects the configuration whose components best address the
diagnosed bottleneck.  The best path is a function of the data, not the system.

Figure~\ref{fig:router} illustrates the full router as a flowchart.  Four orthogonal
axes of the capture profile drive the routing decision: (i)~\textbf{sharpness} --- is
motion blur the dominant defect, or is the footage sharp?; (ii)~\textbf{coverage} ---
are under-observed regions (floaters, holes) the dominant defect, or is the splat dense?;
(iii)~\textbf{capture hardware} --- RGB-only, or does a depth sensor or event camera
complement the video?; and (iv)~\textbf{deliverable} --- is a splat room representation
sufficient, or does the project mandate a polygonal mesh for game-engine collision and
LOD?  The four axes are independent; the router evaluates each and assembles the
appropriate combination of components from the validated menu described in the sections
below.

Critically, the polygonal-mesh mandate can be \emph{decoupled} from the room
representation.  The TRELLIS.2 objects (Section~\ref{sec:method:objects}) are polygonal
meshes (GLB$\rightarrow$FBX) in every viable option; a splat room with polygonal objects
embedded in Unreal Engine~5.8 therefore satisfies the mandate without requiring the room
itself to be meshed.  This decoupling is the key architectural insight that resolves the
conflict between mesh-fidelity requirements and the empirical finding that meshing a
blur-limited room produces output that is not appreciably better than the splat it was
derived from (Section~\ref{sec:results:mesh}).

% ─────────────────────────────────────────────────────────────────────────────
\section{Stage 1: Ingest and Frame-Quality Gate}
\label{sec:method:ingest}

\subsection{Frame extraction}
Video is decoded with PyAV at the native frame rate and stored as lossless PNG to preserve
sub-pixel feature information.  An integrity manifest is written per-capture: each frame
entry carries its source index, full-resolution Laplacian-variance sharpness score, and a
SHA-256 hash of the pixel data; the source video is itself hashed so that any silent
re-extraction can be detected.  The same canonical frame set feeds structure-from-motion
(SfM), Gaussian training, and metrics, ensuring camera-to-image correspondence is
consistent throughout the pipeline.

\subsection{Neural image-quality gate (MUSIQ)}
After extraction, \texttt{frame\_quality.py} scores every frame with the
\textbf{MUSIQ}~\cite{ke2021_musiq} no-reference image-quality assessment model, executed
via \texttt{pyiqa}~\cite{chen2022_pyiqa} on the GPU (classical Laplacian-variance fallback
when the GPU is unavailable).  MUSIQ is perceptually calibrated: on the KonIQ-10k scale
(0--100, higher is better) a photographic still scores 50--75; frames in the teens are
unusably soft.  The gate applies a \emph{drop-and-flag policy}: a video that cannot
supply \texttt{min\_good\_frames} KEEP-grade frames is flagged \texttt{too\_low\_quality}
and skipped so that the orchestrator moves to the next video rather than feeding low-grade
footage into SfM and 3DGS training.  Accepted videos return a \emph{windowed best-of}
selection --- the sharpest frame per overlapping temporal window --- which simultaneously
maximises quality and preserves coverage across the duration of the capture.

\subsection{Directional sharpness gate for motion blur}
\label{sec:method:directional}
The dreamlab reference capture (Section~\ref{sec:results:overall}) revealed a limitation
of scalar Laplacian variance as a blur metric.  At 4K resolution, the per-frame Laplacian
range spans approximately 105$\times$ between the sharpest and blurriest frames, but
down-scaling to even 960p before computing the Laplacian collapses this range to
$\approx$1.6$\times$, making blurry frames appear indistinguishable from sharp ones.
All sharpness assessment is therefore performed at full or near-full resolution.

More fundamentally, scalar Laplacian variance measures the \emph{magnitude} of the second
derivative without regard to spatial orientation.  Uniform motion blur, however, produces
a characteristically \emph{directional} attenuation: edges aligned with the blur direction
lose contrast severely while edges orthogonal to it are less affected.  The upgraded gate
additionally computes the \textbf{structure-tensor minor eigenvalue} $\lambda_2$ for each
frame, which quantifies the degree of isotropic contrast in the gradient field.  A
high-quality frame has abundant gradients in all directions ($\lambda_2$ is large);
a motion-blurred frame has gradients predominantly in one direction ($\lambda_2 \approx
0$).  The gate issues a \textbf{recapture flag} when the population distribution of
$\lambda_2$ across a video is skewed towards near-zero values, signalling that the blur is
systematic rather than isolated to a minority of frames.  Four unit tests validate the
metric against synthetic images with known blur direction and magnitude.

Measurement on the dreamlab capture confirms both metrics agree that the footage is
blur-limited; $\lambda_2$ additionally identifies the dominant blur direction as
translational (horizontal sweeping motion), consistent with the handheld capture
trajectory.

% ─────────────────────────────────────────────────────────────────────────────
\section{Stage 2: Structure from Motion}
\label{sec:method:sfm}

Camera poses are estimated with \textbf{COLMAP}~\cite{schoenberger2016sfm} version 4.1.0.
Feature extraction uses GPU SIFT (the COLMAP 4.1.0 namespace changed to
\texttt{--FeatureExtraction.*}); ALIKED+LightGlue~\cite{lindenberger2023_lightglue} is the
intended SOTA matcher and is wired into the pipeline configuration, though the current
deployment falls back to GPU SIFT because the ALIKED enum is invalid in the COLMAP 4.1.0
build and the LightGlue extension requires \texttt{libcudnn.so.9} (absent from the
container image).  GPU SIFT gave 600/600 frames registered on the dreamlab capture (mean
re-projection error 1.03~px), so the matcher limitation is not the binding constraint on
the current reference run.

For multi-sweep captures where the same space is filmed at multiple heights in separate
passes, all sweeps are ingested together and COLMAP registers them as one unified model.
The multi-sweep strategy demonstrably raises point-cloud density: on a three-sweep
combined capture, COLMAP produced $\sim$146k points at $\sim$95\% registration, versus
$\sim$90k points for a single sweep of the same room.

% ─────────────────────────────────────────────────────────────────────────────
\section{Stage 3: 3D Gaussian Splatting Training}
\label{sec:method:3dgs}

\subsection{The 3DGS representation}
3D Gaussian Splatting~\cite{kerbl2023_3dgs} represents a scene as a set of anisotropic
3D Gaussians, each parameterised by a position $\boldsymbol{\mu} \in \mathbb{R}^3$, a
covariance matrix $\boldsymbol{\Sigma}$ (decomposed as scale $\mathbf{s}$ and rotation
quaternion $\mathbf{q}$), an opacity $\alpha$, and view-dependent colour encoded in
spherical-harmonic coefficients.  Rasterisation sorts Gaussians by depth and
alpha-composites them into the image, yielding real-time rendering rates after training.

\subsection{LichtFeld native trainer}
Vitrine vendors \textbf{LichtFeld Studio}~\cite{lichtfeld2026} (v0.5.2,
C++23/CUDA, MIT licence) as a pinned tool --- a git submodule at
\texttt{vendor/lichtfeld-studio} --- for native 3DGS training and its local MCP control
surface (70+ tools; HTTP endpoint \texttt{:45677}).  The vendored tool is never modified;
updates are applied by bumping the submodule tag.

Training uses the \textbf{ImprovedGS+} (\texttt{igs+}) strategy with the MRNF
densification schedule and \textbf{appearance conditioning} (\texttt{--ppisp},
per-camera photo-image signal processing) for correcting per-frame exposure drift.  The
splat is trained \emph{unpruned}: opacity-sparsity regularisation is deliberately
disabled to preserve fidelity for downstream mesh extraction and visual inspection.
Performance in the game engine is handled separately via Nanite (polygonal assets) and
the NanoGS plugin (splat rendering), so the training budget is not constrained by
real-time rendering.  Training to 30,000 iterations produces the canonical
\texttt{splat\_30000.ply} used downstream.  On the dreamlab capture, 2.48M Gaussians were
trained; a floater-prune pass (SuperSplat) reduced this to 1.55M before UE ingest.

% ─────────────────────────────────────────────────────────────────────────────
\section{Stage 4: Scene Representation}
\label{sec:method:scene}

The router selects one of three scene representations depending on the capture profile and
deliverable requirement (Figure~\ref{fig:router}, ``Room deliverable?'' gate).

\subsection{Scene as splat: LichtFeld .ply → SuperSplat → NanoGS}
\label{sec:method:scene:splat}
The primary spine for the dreamlab capture --- and the recommended default for any
blur-limited or coverage-dense RGB capture --- is to \textbf{preserve the splat as the
room representation} and render it as real Gaussians in UE~5.8 via the
\textbf{NanoGS}~\cite{nanogs2026} plugin.

The process is: (i)~export the LichtFeld-trained \texttt{splat\_30000.ply}; (ii)~clean
floaters and crop the field boundary in \textbf{SuperSplat}\footnote{SuperSplat, PlayCanvas, MIT licence. \url{https://github.com/playcanvas/supersplat}} (MIT);
(iii)~import into UE~5.8 via the NanoGS plugin, which was recompiled for UE~5.8 from its
UE~5.6/5.7 codebase.  NanoGS renders real Gaussians via a Nanite-style LOD pipeline,
replacing the LiDAR-point-cloud workaround used prior to this work.  The full dreamlab
scene (pruned splat room + chair, vacuum cleaner, and toolbox FBXs) stands in the editor
and renders correctly (Figure~\ref{fig:nanogs_room}, Figure~\ref{fig:scene_integrated}).

This approach has four properties that make it the preferred path: it uses the
LichtFeld-native trainer output directly (lean-on-vendor principle); it delivers the
highest fidelity available from the captured data without the lossy geometry-extraction
step; it eliminates the flaky UE Remote-Control assembly scripting used in earlier
versions; and it satisfies the polygonal-mesh mandate through the object FBXs rather than
requiring the room to be meshed.

\subsection{Scene as mesh: 2DGS / PGSR / TSDF → FBX (Nanite)}
\label{sec:method:scene:mesh}
When a polygonal room mesh is explicitly required --- for collision volumes, physics
simulation, or a deliverable format that cannot accommodate a splat plugin --- the router
selects the mesh branch.

\textbf{2D Gaussian Splatting}~\cite{huang2024_2dgs} (\textbf{2DGS}) trains
planar surfels rather than 3D ellipsoids; the surfel depth maps are geometrically cleaner
than ellipsoid-depth (which smears across surface boundaries) and produce improved TSDF
fusion input.  \textbf{PGSR}~\cite{chen2024_pgsr} uses planar Gaussians with multi-view
consistency constraints and an explicit textureless-wall prior, making it well-suited to
indoor environments with large featureless surfaces.  Both are kept in the toolkit for
sharp captures; their benefit is empirically absent on the blur-limited dreamlab data
(Section~\ref{sec:results:mesh}).

The TSDF fallback (\texttt{gsplat}-to-\texttt{ScalableTSDFVolume}) is stable for all
capture qualities and remains the default mesh backend.  The Open3D tensor-API
\texttt{VoxelBlockGrid} extractor was evaluated and found to abort with an integer
overflow at scale (an Open3D~0.19 bug); the legacy CPU \texttt{ScalableTSDFVolume} path
is the validated alternative.

Irrespective of the extraction method, the mesh path follows a fixed cleaning and baking
recipe: (i)~\texttt{MeshCleaner.clean(smooth\_iterations=0)} (smoothing introduces
degenerate triangles that crash \texttt{xatlas}); (ii)~\texttt{xatlas} UV-atlas
generation and albedo bake from vertex colours (\texttt{bake\_from\_vertex\_colors});
(iii)~gravity/scale conversion to UE centimetre units; (iv)~FBX export via Blender
(\texttt{blender\_obj\_to\_fbx}); (v)~UE import with \texttt{import\_materials=True}
and Nanite enabled.  Unreal Engine discards vertex colours and the USD
\texttt{displayColor} primvar on import; a baked UV texture is therefore the only
mechanism by which captured colour survives into the engine.

\subsection{Scene as both: splat hero + mesh asset (dreamlab choice)}
For the dreamlab reference run, both representations are produced: the splat serves as
the high-fidelity visual hero (rendered via NanoGS, Figure~\ref{fig:nanogs_room}), while
the TSDF-extracted FBX provides a polygonal asset for the mesh-asset archive
(Section~\ref{sec:results:mesh}).  This dual-representation strategy hedges delivery
format and provides the comparison data reported in Section~\ref{sec:results:mesh}.

% ─────────────────────────────────────────────────────────────────────────────
\section{Stage 5: Object Reconstruction}
\label{sec:method:objects}

The object track runs in parallel with the scene track and is mandatory in every router
configuration; TRELLIS.2 objects satisfy the polygonal-mesh mandate even when the room is
delivered as a splat.

\subsection{Segmentation: SAM3}
Concept-based segmentation is performed with \textbf{SAM3}~\cite{sam3dobjects2025}, which
extends the Segment Anything paradigm to a vocabulary of 4M concepts.  SAM3 is run
per-frame on a sparse temporal sample of the training frames; per-frame masks are unioned
by concept label to produce a per-object binary mask volume.  A depth-aware multi-view
projection then votes each Gaussian centre inside or outside the mask across all
registered cameras that sampled it; Gaussians inside the mask in a majority of views are
assigned to the corresponding object.  This produces per-object Gaussian subsets ranked
by curatorial significance (keyness $\times$ Gaussian count).

\subsection{TRELLIS.2 hull\_e2e}
\label{sec:method:trellis}
Each per-object Gaussian subset is fed to the \textbf{TRELLIS.2}~\cite{xiang2025_trellis2}
\texttt{hull\_e2e} reconstruction path.  TRELLIS.2 is a 4B-parameter multi-view 3D
generation model (MIT licence; 17--60~s per object on a single 48~GB GPU at
\texttt{1536\_cascade} geometry / 4096~px PBR textures).

The hull pipeline is: (i)~orbit-render 6 panels from the per-object PLY (front, left,
back, right at elevation 0$^\circ$; top at $+85^\circ$; bottom at $-85^\circ$);
(ii)~per-panel coverage gate --- panels with alpha fraction below \texttt{gap\_threshold}
are synthesised by \textbf{FLUX.2}~\cite{flux2025} view-completion conditioned on
observed panels, recovering geometry for faces the camera never reached;
(iii)~\texttt{Trellis2MultiViewImageToShape} shape diffusion on all six panels;
(iv)~\texttt{Trellis2RasterizePBR} producing a textured GLB with base colour,
metallic/roughness, and alpha maps at 4096~px.

The GLB is pre-scaled to real-world centimetres via \texttt{blender\_prescale\_objects}
(reading physical dimensions from the COLMAP point cloud), then exported to FBX for UE
import.  The TRELLIS.2 FBX assets import into UE~5.8 with \texttt{import\_materials=True}
and Nanite enabled; object placement uses \texttt{ue\_place\_prescaled} with
pre-computed real-world coordinates to avoid the flaky \texttt{get\_actor\_bounds}
dependency.\footnote{The \texttt{get\_actor\_bounds} Remote Control call was observed to throw on final UE assembly; pre-computed coordinates from \texttt{object\_meta.json} sidestep the call entirely.}

TRELLIS.2 is empirically the superior object reconstruction method over Hunyuan3D-2.1 for
the dreamlab objects (chair, vacuum cleaner, toolbox): its hull outputs are cleaner and
more recognisable given the same per-object PLY input.  Hunyuan3D-2.1~\cite{zhao2025_hunyuan3d20}
is retained as a fallback for cluttered or partially-observed objects.

% ─────────────────────────────────────────────────────────────────────────────
\section{Stage 6: Capture-Conditional Enhancers}
\label{sec:method:enhancers}

The router attaches capture-conditional enhancement stages whose benefit is specific to
a diagnosed bottleneck.  These are not part of the default spine.

\subsection{ArtiFixer: floater and under-observation repair}
\label{sec:method:artifixer}
\textbf{ArtiFixer}~\cite{delutio2026_artifixer} is a 14B-parameter 3DGS reconstruction-enhancement
model (NVIDIA, Apache-2.0 code; weight licence: NVIDIA OneWay Noncommercial) that
operates in the Gaussian parameter space to regenerate \emph{under-observed} regions ---
holes and floaters caused by insufficient view coverage --- by learning to fill them from
observed context.  ArtiFixer addresses a distinct failure mode from motion blur: it is
effective when the scene is under-observed (coverage gaps), and is explicitly not a
deblurring tool.

The router activates ArtiFixer only on captures diagnosed as \emph{coverage-limited}
(holes, floaters, $>0\%$ under-observed area) and explicitly deactivates it on
captures diagnosed as \emph{blur-limited}, where its output is wasted compute.  The
dreamlab capture has 0\% under-observed area (median $\sim$90 cameras per surface point)
and a MUSIQ score of $\sim$31 (blur-limited), so ArtiFixer is \textbf{not} applied on
this run.

ArtiFixer runs in an optional sidecar container (\texttt{gaussian-toolkit-artifixer:ada-sm89})
on a dedicated GPU.  Section~\ref{sec:results:artifixer} reports the OOM analysis and
hardware characterisation performed on the Ada RTX~6000 (sm\_89).

\subsection{BAD-Gaussians / 3dgs-deblur (experimental)}
\label{sec:method:deblur}
\emph{Faithful} joint deblur-during-reconstruction --- methods that model the exposure
trajectory as a learnable parameter and jointly optimise camera motion and 3DGS
parameters~\cite{zhao2024_badgaussians} --- is qualitatively distinct from generative
per-frame pre-deblur (NAFNet/Restormer), which is rejected because it hallucinates
fine detail that breaks multi-view consistency.  The faithful joint approach is the one
honest non-recapture lever for blur-limited rooms; it is wired into the component menu
(Figure~\ref{fig:router}, ``S\_BLUR'' branch) but has not been run on the dreamlab data
as of this report.  Running a falsification test is the highest-priority experiment for
the blur-limited room track.

% ─────────────────────────────────────────────────────────────────────────────
\section{Stage 7: Unreal Engine 5.8 Assembly}
\label{sec:method:ue}

Delivery targets \textbf{Unreal Engine~5.8}, which shipped June 2026 with a first-party
experimental MCP plugin.  The UE~5.8 overlay runs in an optional Docker container
(\texttt{vitrine-unreal:5.8}), built from an NVIDIA CUDA base with the UE Linux runtime
installed in-repo at \texttt{unreal/engine/} ($\approx$73~GB, bind-mounted read-only).

The control path is the UE Web Remote Control API (port~30010, unicast), proxied through
\texttt{unreal-mcp-bridge} (port~9100) on the shared Docker network.  The first-party
UE MCP (port~8000) is available as a secondary path.

Assembly follows a pre-scale-offline strategy to avoid the flaky
\texttt{get\_actor\_bounds} Remote Control call: object FBXs are scaled to real-world
centimetres before import using \texttt{blender\_prescale\_objects}; placement coordinates
are computed from \texttt{object\_meta.json} (offline COLMAP measurements).
\texttt{ue\_place\_prescaled} then imports and places each asset at \texttt{scale=1.0}
with pre-computed coordinates, no actor-query API calls required.  The room FBX is
imported with Nanite enabled; the splat (when used as the room hero) is imported via the
NanoGS plugin as described in Section~\ref{sec:method:scene:splat}.
